% !TeX root = ../Cosmology.tex

%% --------------------------------------------------------
\chapter{Поляризація CMB}
\label{sec:pol}
%% --------------------------------------------------------

Температурний спектр $C_\ell^{TT}$, розглянутий у попередньому
розділі, несе інформацію переважно про скалярні збурення.
Але CMB містить і другий канал інформації~--- \emph{поляризацію}.
Вона виникає при томсонівському розсіянні фотонів і дозволяє
розрізнити два типи первинних збурень: скалярні (флуктуації
густини) і тензорні (гравітаційні хвилі). Детекція тензорного
сигналу у поляризації CMB стала б прямим підтвердженням
квантової природи гравітації~--- саме тому пошук B-мод є
однією з центральних задач сучасної космології.

%% --------------------------------------------------------
\section{Механізм поляризації: розсіяння Томсона}
%% --------------------------------------------------------

Фотони поляризуються при останньому розсіянні на електронах.
Томсонівське розсіяння поляризує фотон \emph{лише} якщо навколо електрона
існує \textbf{квадрупольна} температурна анізотропія: якщо температура
однакова з усіх боків~--- поляризації немає.

Безпосередньо до рекомбінації фотонів і електронів ще достатньо для
ізотропізації, тому квадруполь малий і поляризація слабка.
Але під час самої рекомбінації ($\Delta z \sim 80$) оптична глибина
падає швидко~--- квадруполь встигає вирости і поляризація «відбивається»
в поверхні останнього розсіяння.

Напрямок поляризації залежить від \textbf{поперечного} перепаду температури
навколо точки розсіяння.

%% --------------------------------------------------------
\section{Параметри Стокса та спін-2 характер поляризації}
\label{subsec:stokes}
%% --------------------------------------------------------

Лінійна поляризація у заданому напрямку $\hat{\mathbf{n}}$ на небесній
сфері описується двома параметрами Стокса $Q(\hat{\mathbf{n}})$ та
$U(\hat{\mathbf{n}})$, визначеними відносно локального базису
$(\hat{\boldsymbol\theta}, \hat{\boldsymbol\phi})$: $Q$ вимірює різницю
інтенсивностей уздовж цих двох осей, а $U$~--- те саме, але для осей,
повернутих на $45^\circ$. Третій параметр Стокса, $V$ (циркулярна
поляризація), точно дорівнює нулю для томсонівського розсіяння
неполяризованого випромінювання~--- це чисто лінійний ефект.

Ключова відмінність $Q,U$ від звичайного скалярного поля (як $\Theta_0$)
у тому, що вони \textbf{не інваріантні} відносно повороту локального
базису. Якщо повернути осі на кут $\psi$, параметри перетворюються як
\begin{equation}\label{eq:QU_rotation}
    (Q \pm iU)'(\hat{\mathbf{n}}) = e^{\mp 2i\psi}\,(Q \pm iU)(\hat{\mathbf{n}}).
\end{equation}
Множник $e^{\mp 2i\psi}$ (а не $e^{\mp i\psi}$, як для вектора) означає,
що комбінація $Q\pm iU$ є полем \textbf{спіральності (спін-ваги) $\pm2$}:
поворот базису на $180^\circ$ повертає поляризацію саму на себе~---
це і є те «безголове» (headless) означення напрямку, використане на
рис.~\ref{fig:eb-modes}.

Оскільки $Q\pm iU$ не є звичайними скалярами, їх не можна розкласти за
звичайними сферичними гармоніками $Y_{\ell m}$. Натомість використовують
\emph{спін-зважені сферичні гармоніки} $\,_{\pm2}Y_{\ell m}(\hat{\mathbf{n}})$
\cite{zaldarriaga_seljak_1997, kamionkowski_kosowsky_stebbins_1997}:
\begin{equation}\label{eq:QU_expansion}
    (Q \pm iU)(\hat{\mathbf{n}}) = \sum_{\ell m} a_{\pm 2,\,\ell m}\;\, _{\pm2}Y_{\ell m}(\hat{\mathbf{n}}).
\end{equation}

%% --------------------------------------------------------
\section{E- і B-моди}
%% --------------------------------------------------------

Поляризаційне поле на сфері, як і будь-яке тензорне поле,
однозначно розкладається на дві компоненти~--- аналогічно
до розкладу вектора на потенціальну і вихрову частини.
Ці дві компоненти мають різне фізичне джерело і різний
статус спостережень (рис.~\ref{fig:eb-modes}):

%% --------------------------------------------------------
\begin{figure}[htbp]
\centering
\begin{tikzpicture}[
    line cap=round,
    every node/.style={font=\footnotesize},
]

\definecolor{Emode}{RGB}{200,40,40}
\definecolor{Bmode}{RGB}{40,60,180}

% ---------- macro: two-ring polarization pattern with a soft central spot ----------
% #1 = x-shift, #2 = y-shift, #3 = colour, #4 = angle offset
%   (0 = radial -> E<0 ; 90 = tangential -> E>0 ; +45/-45 = B<0/B>0)
% #5 = caption text (placed BELOW the panel, TikZ math mode allowed)
\newcommand{\polfield}[5]{%
    \begin{scope}[xshift=#1cm, yshift=#2cm]
        % soft central spot -- just a hint of "there's a perturbation here",
        % not tied to a specific hot/cold sign
        \shade[ball color=#3!12, opacity=0.55] (0,0) circle (0.42);
        % inner ring
        \foreach \ang in {0,45,...,315}{
            \draw[#3, line width=1.05pt]
                ({0.28*cos(\ang+#4)+0.62*cos(\ang)},{0.28*sin(\ang+#4)+0.62*sin(\ang)})
                --
                ({-0.28*cos(\ang+#4)+0.62*cos(\ang)},{-0.28*sin(\ang+#4)+0.62*sin(\ang)});
        }
        % outer ring, staggered by 22.5 deg, slightly longer sticks
        \foreach \ang in {22.5,67.5,...,337.5}{
            \draw[#3, line width=1.05pt]
                ({0.34*cos(\ang+#4)+1.15*cos(\ang)},{0.34*sin(\ang+#4)+1.15*sin(\ang)})
                --
                ({-0.34*cos(\ang+#4)+1.15*cos(\ang)},{-0.34*sin(\ang+#4)+1.15*sin(\ang)});
        }
        % caption below the panel, outside the pattern
        \node at (0,-1.75) {#5};
    \end{scope}
}

\polfield{0}{4.2}{Bmode}{0}{$E<0$}
\polfield{6}{4.2}{Emode}{90}{$E>0$}
\polfield{0}{0}{Bmode}{45}{$B<0$}
\polfield{6}{0}{Emode}{-45}{$B>0$}

\end{tikzpicture}
\caption{Характерні візерунки поляризації навколо точки останнього
розсіяння. Штрихи позначають напрямок поляризації і навмисно
без стрілок~--- поляризація є \emph{безголовим} напрямком (визначена
з точністю до $180^\circ$, а не $360^\circ$, як звичайний вектор).
\textbf{E-моди} ($E<0$, $E>0$)~--- радіальний або тангенціальний
візерунок, парний відносно віддзеркалення (як і скалярне джерело,
що його породжує). \textbf{B-моди} ($B<0$, $B>0$)~--- візерунок,
повернутий на $45^\circ$ відносно E-моди, з вираженою закрученістю
(handedness), що й відповідає їхній непарній паритетній природі.}
\label{fig:eb-modes}
\end{figure}
%% --------------------------------------------------------

\begin{center}
\begin{tblr}{XXX}
\toprule
Тип & Джерело & Статус \\
\midrule
E-моди (градієнтний, симетричний візерунок) & Скалярні флуктуації & Виміряні \\
B-моди (ротаційний, закручений візерунок)  & Тензорні флуктуації (GW) & Поки не знайдені \\
\bottomrule
\end{tblr}
\end{center}

Чому два типи збурень дають різні моди поляризації~---
питання симетрії, а не лише статистики:

\begin{keybox}{Чому скалярні збурення не дають B-мод?}
Скалярні збурення мають потенціальний (паритетно-парний) характер:
$\Phi(\mathbf{x})$ визначає поле швидкостей $\mathbf{v} = \nabla\Phi$.
E-моди мають ту ж симетрію паритету~--- вони парні.
B-моди псевдоскалярні (непарні), і лінійне скалярне збурення їх не генерує.
Тензорне збурення $h_{ij}$ має обидва стани, в тому числі непарний~---
звідси B-моди від гравітаційних хвиль.
\textit{Практичний забруднювач:} гравітаційне лінзування на великих масштабах перетворює
E-моди в B-моди~--- це треба враховувати при пошуку первинних B-мод.
\end{keybox}

Формально E- і B-моди означуються через (анти)симетричну комбінацію
коефіцієнтів $a_{\pm2,\ell m}$ з рівняння~\eqref{eq:QU_expansion}:
\begin{equation}\label{eq:E_B_def}
    a_{\ell m}^{E} = -\tfrac{1}{2}\left(a_{2,\ell m} + a_{-2,\ell m}\right),
    \qquad
    a_{\ell m}^{B} = \tfrac{i}{2}\left(a_{2,\ell m} - a_{-2,\ell m}\right).
\end{equation}
Під дією просторової інверсії $\hat{\mathbf{n}} \to -\hat{\mathbf{n}}$
скалярні гармоніки $Y_{\ell m}$ перетворюються як $(-1)^\ell$, тоді як
$a_{\ell m}^E$ успадковує цю ж парність $(-1)^\ell$ (як і температура
$a_{\ell m}^T$), а $a_{\ell m}^B$ отримує \emph{додатковий} знак~---
парність $(-1)^{\ell+1}$. Саме ця різниця у трансформаційних
властивостях, а не якась відмінність в амплітудах, і є строгим
математичним змістом твердження в рамці вище.

%% --------------------------------------------------------
\section{Кутові спектри поляризації}
\label{subsec:pol_spectra}
%% --------------------------------------------------------

За аналогією з $C_\ell^{TT}$ означують кутові спектри потужності
для всіх пар полів $X,Y \in \{T,E,B\}$:
\begin{equation}\label{eq:pol_spectra_def}
    \langle a_{\ell m}^{X}\, a_{\ell' m'}^{Y*} \rangle
    = C_\ell^{XY}\, \delta_{\ell\ell'}\, \delta_{mm'}.
\end{equation}
Якщо первинні збурення (скалярні й тензорні) паритетно-симетричні,
як у стандартній інфляційній картині, то крос-спектри $C_\ell^{TB}$ і
$C_\ell^{EB}$ мусять тотожно дорівнювати нулю~--- вони змінюють знак
під інверсією, а статистично ізотропний паритетно-симетричний
Всесвіт не може виділити такий знак. Ненульові $C_\ell^{TB}$ або
$C_\ell^{EB}$ є або систематичною помилкою вимірювання (найчастіше),
або, у принципі, ознакою нової паритетно-непарної фізики.

Натомість $C_\ell^{TE}$ дозволений і несе важливу інформацію: обидва
поля, $T$ і $E$, походять з одного джерела~--- акустичних осциляцій
плазми до рекомбінації. Але поляризацію сорсує не сама монопольна
температура $\Theta_0$, а швидкість баріонів $v_b \propto \dot\Theta_0$,
яка зсунута за фазою на $\pi/2$ відносно $\Theta_0$ у акустичному
осциляторі. Це дає характерну картину:
\begin{itemize}
    \item Спектр $C_\ell^{EE}$ також осцилює з піками акустичної
          природи, але його піки \textbf{зсунуті за фазою} відносно
          піків $C_\ell^{TT}$~--- там, де температурний спектр має
          екстремум, поляризаційний проходить через нуль (і навпаки).
    \item Спектр $C_\ell^{TE}$ змінює знак кілька разів; перший
          виражений провал (антикореляція) на $\ell \sim 150$ є
          прямим наслідком \textbf{адіабатичності} первинних збурень
          і використовується як незалежний тест проти ізокривизняних
          альтернатив.
    \item При малих мультиполях ($\ell \lesssim 10$) у $C_\ell^{EE}$
          з'являється додатковий \textbf{бамп реіонізації}: фотони
          повторно розсіюються на вільних електронах під час епохи
          реіонізації ($z_{\text{re}} \sim 7$--$10$), генеруючи новий
          квадруполь і, отже, нову поляризацію на надгоризонтних
          масштабах цієї епохи. Амплітуда цього бампу~--- практично
          єдиний спосіб незалежно виміряти оптичну глибину реіонізації
          $\tau$, не покладаючись на форму плато Закса-Вольфа в $TT$.
\end{itemize}

\begin{figure}[htbp]
\centering
\begin{tikzpicture}
\begin{axis}[
    width=0.95\textwidth,
    height=6.5cm,
    xlabel={Мультиполь $\ell$},
    ylabel={$\mathcal{D}_\ell$ [умовні одиниці]},
    xmin=2, xmax=1000,
    xmode=log,
    ymin=-2, ymax=6,
    legend pos=north east,
    legend style={font=\scriptsize},
    axis lines=left,
    grid=major,
    grid style={gray!15},
]
\addplot[
    domain=2:1000, samples=300, thick, color=orange!85!black,
] {2.2*exp(-x/900)*(1+0.7*cos(3.14159*(x)/165))^2 + (x<15)*1.1*exp(-(x-2)/4)};
\addlegendentry{$C_\ell^{EE}\cdot \ell^2$ (схематично)}

\addplot[
    domain=2:1000, samples=300, thick, dashed, color=teal,
] {1.8*exp(-x/700)*sin(3.14159*(x-40)/165)};
\addlegendentry{$C_\ell^{TE}\cdot \ell^2$ (схематично)}

\draw[->, gray!70, thin] (axis cs:5,3.6) -- (axis cs:9,1.9);
\node[font=\tiny, text=gray!70, align=center] at (axis cs:6,4.1) {бамп\\реіонізації};
\end{axis}
\end{tikzpicture}
\caption{Якісна ілюстрація форми спектрів $C_\ell^{EE}$ та
$C_\ell^{TE}$ (не апроксимація реальних даних Planck, лише для
пояснення фазових співвідношень). Зверніть увагу на зсув фази піків
$EE$ відносно $TT$ (рис.~\ref{fig:cmb_tt} у попередньому розділі),
знакозмінність $TE$ та додатковий бамп реіонізації при
$\ell \lesssim 10$.}
\label{fig:ee_te_spectra}
\end{figure}

Спектр $C_\ell^{BB}$ на цьому ж малюнку показати важко: первинний
(тензорний) сигнал за амплітудою на $2$--$3$ порядки менший за
$C_\ell^{EE}$, і додатково існує \emph{вторинний} $C_\ell^{BB}$,
породжений гравітаційним лінзуванням E-мод великомасштабною
структурою (див. далі).

%% --------------------------------------------------------
\section{Тензорно-скалярне відношення $r$}
\label{subsec:tensor_scalar_ratio}
%% --------------------------------------------------------

Амплітуда первинних B-мод параметризується \emph{тензорно-скалярним
відношенням}
\begin{equation}\label{eq:r_def}
    r \equiv \frac{A_t}{A_s},
\end{equation}
де $A_t$ і $A_s$~--- амплітуди спектрів потужності тензорних (гравітаційні
хвилі) і скалярних (густина) збурень у момент виходу за горизонт під
час інфляції. Для однопольових моделей повільного кочення це відношення
безпосередньо пов'язане з енергетичним масштабом інфляції:
\begin{equation}\label{eq:r_energy_scale}
    V^{1/4} \approx 1.06\times10^{16}\,\text{ГеВ}\times\left(\frac{r}{0.01}\right)^{1/4}.
\end{equation}
Додатково, у цьому класі моделей діє \emph{consistency relation}~---
співвідношення, яке пов'язує $r$ з нахилом тензорного спектра $n_t$
без жодного вільного параметра:
\begin{equation}\label{eq:consistency_relation}
    n_t = -\frac{r}{8}.
\end{equation}
Перевірка цього співвідношення (потребує вимірювання й $r$, і $n_t$
окремо) є одним із найчистіших тестів однопольової інфляції~---
але при поточній чутливості вимірити $n_t$ окремо від $r$ поки
нереалістично.

%% --------------------------------------------------------
\section{Практичні виклики: фон, лінзування і delensing}
\label{subsec:foregrounds}
%% --------------------------------------------------------

Крім гравітаційного лінзування, згаданого в рамці вище, пошук
первинних B-мод ускладнюють два додаткові ефекти.

\textbf{Галактичні форграунди.} Поляризоване випромінювання пилу
(теплового) і синхротронне випромінювання від Чумацького Шляху
самі створюють B-моди на тих самих кутових масштабах, що й
первинний сигнал, причому за амплітудою можуть суттєво його
перевищувати. Розділити космологічний сигнал і форграунд можна лише
завдяки різній частотній залежності~--- звідси потреба у
багаточастотних вимірюваннях (десятки--сотні ГГц).

\begin{keybox}{Урок BICEP2 (2014)}
У березні 2014 року колаборація BICEP2 оголосила про детекцію
первинних B-мод з $r \approx 0.2$~--- сенсаційний результат, що
відповідав би прямому спостереженню гравітаційних хвиль з інфляції.
Подальший спільний аналіз з даними Planck по пилу показав, що
основну частину сигналу пояснює поляризоване випромінювання
галактичного пилу, яке недооцінили на момент першої публікації
\cite{bicep2_2014, planck_bicep2_joint_2015}. Цей епізод~--- стандартний
приклад того, чому детекція B-мод вимагає впевненого контролю
форграундів на кількох частотах, перш ніж заявляти про
космологічне походження сигналу.
\end{keybox}

\textbf{Delensing.} Оскільки лінзові B-моди породжуються відомим
(і вимірюваним незалежно) полем E-мод та гравітаційним потенціалом
великомасштабної структури, їхній внесок можна статистично
відняти~--- процедура, відома як \emph{delensing}. Це один із
головних напрямів розвитку наступного покоління експериментів,
оскільки без нього лінзовий B-сигнал стає ефективною підлогою
чутливості до $r$ на рівні $r \sim 10^{-3}$.

%% --------------------------------------------------------
\section{Огляд поточних і майбутніх експериментів}
\label{subsec:experiments}
%% --------------------------------------------------------

\begin{center}
\begin{tblr}{XXXX}
\toprule
Експеримент & Тип & Стан & Ціль по $r$ \\
\midrule
Planck & Супутник & Завершено (2018) & --- \\
BICEP/Keck & Наземний & Діючий & $r < 0.036$ \cite{bicep_keck_2021} \\
Simons Observatory & Наземний & Будується & $\sigma(r) \sim 0.002$--$0.003$ \\
CMB-S4 & Наземний & Планується & $\sigma(r) \sim 0.0005$ \\
LiteBIRD & Супутник & Планується (2030-ті) & $\sigma(r) \sim 0.001$ \\
\bottomrule
\end{tblr}
\end{center}

Знайти B-моди~--- означає отримати прямий доказ первинних гравітаційних хвиль
і квантової природи гравітації в режимі малих збурень.
Поточна верхня межа: $r < 0.036$ (95\% CL, BICEP/Keck~2021 \cite{bicep_keck_2021}).

Таким чином, CMB є повноцінним «двоканальним» джерелом інформації:
температурний спектр обмежує скалярні збурення і параметри $\Lambda$CDM,
а поляризаційні B-моди~--- єдиний прямий зонд тензорних збурень і
енергетичного масштабу інфляції. Обидва канали разом утворюють
найпотужніший інструмент для перевірки інфляційних моделей~---
включно з тією, що розглядається у наступних розділах.

%% --------------------------------------------------------
\section*{Задачі}
\addcontentsline{toc}{section}{Задачі}
%% --------------------------------------------------------

\begin{problem}{Спін-2 характер поляризації}
Розглянемо поворот локального базису на кут $\psi$, при якому
параметри Стокса перетворюються за формулою~\eqref{eq:QU_rotation}.
\begin{enumerate}
    \item Покажіть, що звичайна температура $\Theta_0$ (спін $0$)
          не змінюється при такому повороті, тоді як $Q+iU$ отримує
          фазовий множник $e^{-2i\psi}$.
    \item Поясніть, чому саме множник $e^{\mp 2i\psi}$, а не
          $e^{\mp i\psi}$, унеможливлює зображення поляризації
          звичайною стрілкою (вектором).
    \item Чому при $\psi = 180^\circ$ поляризаційний «штрих» переходить
          сам у себе, а не у протилежний напрямок?
\end{enumerate}
\end{problem}

% -------------------------------------------------------------------

\begin{problem}{Парність $E$- і $B$-мод}
Використовуючи означення~\eqref{eq:E_B_def} та властивість парності
звичайних сферичних гармонік $Y_{\ell m}(-\hat{\mathbf n}) = (-1)^\ell
Y_{\ell m}(\hat{\mathbf n})$:
\begin{enumerate}
    \item Обґрунтуйте, чому $a_{\ell m}^E$ має парність $(-1)^\ell$
          (така сама, як у температури $a_{\ell m}^T$).
    \item Обґрунтуйте, чому $a_{\ell m}^B$ отримує додатковий знак
          і має парність $(-1)^{\ell+1}$.
    \item Використовуючи ці результати, покажіть, чому статистично
          ізотропний і паритетно-симетричний Всесвіт вимагає
          $C_\ell^{TB} = C_\ell^{EB} = 0$.
\end{enumerate}
\end{problem}

% -------------------------------------------------------------------

\begin{problem}{Тензорно-скалярне відношення та енергетичний масштаб інфляції}
Скористайтесь співвідношенням~\eqref{eq:r_energy_scale}.
\begin{enumerate}
    \item Оцініть енергетичний масштаб інфляції $V^{1/4}$ для
          поточної верхньої межі $r < 0.036$.
    \item У скільки разів має покращитись чутливість до $r$
          (порівняно з поточною межею), щоб знизити оцінку $V^{1/4}$
          вдвічі?
    \item Скориставшись consistency relation~\eqref{eq:consistency_relation},
          оцініть $n_t$ для $r = 0.01$. Чи реалістично виміряти таке
          значення $n_t$ окремо найближчим часом? Обґрунтуйте відповідь,
          користуючись таблицею цілей експериментів вище.
\end{enumerate}
\end{problem}

% -------------------------------------------------------------------

\begin{problem}{Оптична глибина реіонізації з бампу EE}
Бамп реіонізації в $C_\ell^{EE}$ на масштабах $\ell \lesssim 10$
за амплітудою приблизно пропорційний $\tau^2$ (у наближенні малих
$\tau$), тоді як внесок первинного плато Закса-Вольфа в $C_\ell^{TT}$
на тих самих масштабах практично не залежить від $\tau$.
\begin{enumerate}
    \item Поясніть фізично, чому амплітуда $EE$-бампу залежить саме
          від $\tau^2$, а не від $\tau$ лінійно (підказка: скільки
          разів фотон повинен розсіятися, щоб набути поляризації
          \emph{і} щоб цей ефект дійшов до спостерігача з правильною
          статистичною вагою).
    \item Чому вимірювання $\tau$ саме через поляризацію (а не лише
          через форму $C_\ell^{TT}$) істотно зменшує вироджену
          залежність між $\tau$ та амплітудою скалярних збурень $A_s$
          у $\Lambda$CDM-підгонці?
    \item Якою мірою помилка у вимірюванні $\tau$ впливає на оцінку
          $r$, отриману з низько-$\ell$ частини спектра $C_\ell^{BB}$?
\end{enumerate}
\end{problem}
